% Source: Solutions/Karan-Gangadharan-PS7-Q1-4.tex (LaTeX source). A stray "\]" after the boxed answer was dropped.
\submittedby{Karan Gangadharan}

The quotient topology on $A$ induced by $p$ is given by
\[ \tau_A=\{U\subseteq A : p^{-1}(U)\text{ is open in }\R\}. \]
Since
% REVIEW: the labels are swapped relative to Munkres's Example 3, where p(x)=a for x>0 and
% p(x)=b for x<0; the student never defines p. The resulting topology is the same by symmetry.
\[ p^{-1}(\{a\})=(-\infty,0),\qquad p^{-1}(\{b\})=(0,\infty),\qquad p^{-1}(\{c\})=\{0\}, \]
and
\[ p^{-1}(\{a,b\})=\R\setminus\{0\}, \]
we observe that the open subsets of $A$ are
\[ \varnothing,\quad\{a\},\quad\{b\},\quad\{a,b\},\quad\{a,b,c\}. \]
Hence, the quotient topology on $A$ is
\[ \boxed{\tau_A=\{\varnothing,\{a\},\{b\},\{a,b\},\{a,b,c\}\}.} \]
